\section{Conclusion}

This paper recast MANAR as an evidence-scoped pre-flight software testbed rather than a flight-ready rescue system. The current artifact represents Haversine navigation, nearest-neighbor sector ordering, manual-order preservation, plan locking, RTL state changes, and ONNX detector harnesses in source. A fixed-seed 500-instance route study shows that nearest-neighbor tours exceed a 2-opt-refined local-search reference by $4.63\%$--$11.09\%$ on average across the evaluated target counts. A separate 26-clip diagnostic records lower model-session latency and wider positive-scenario coverage for YOLO11n than D-FINE-N, but the absence of bounding-box ground truth restricts interpretation to raw events and clip coverage. The most important result is therefore the boundary itself: the repository supports a reproducible software baseline and several design hypotheses, not physical endurance, detector accuracy, multisensor fusion, or operational rescue performance. Closing that gap requires annotations, broader baselines, automated tests, end-to-end target-hardware measurements, SITL/HITL, and controlled flight evidence.
